\section{Step 16: Separate full-frame shape from delta shape}
A full frame has exactly 4096 RGB triples, or 12\,288 bytes. Its length is
sufficient to prove the shape because every byte value is an allowed channel.
A delta uses variable-length commands, so it needs a cursor walk before it can
be applied.

Fill \texttt{record\_shape}. For kind zero, require the saved payload length to
be 12\,288. For both kinds, initialize the payload cursor \texttt{\$t0} from
\texttt{\$s5} and the pixel cursor \texttt{\$t1} to zero. Supplied control
flow sends a validated full frame directly to the application boundary. A delta
enters command validation. Neither route has yet changed a pixel.

\paragraph{Observe} Prepare \texttt{s16-1} and stop before
\texttt{record\_shape\_ready}. At \texttt{record\_shape\_ready}, the worked delta has
\texttt{\$t0} equal to source plus sixteen and \texttt{\$t1} zero.

\Needspace{7\baselineskip}
\paragraph{Check} Why does a full frame not need a command walk?
\begin{itemize}
\item[A.] Its triples have fixed extent and no command headers.
\item[B.] Its CRC is always zero.
\item[C.] Its old pixels supply missing bytes.
\end{itemize}
